\section{Revisão bibliográfica}
\label{sc:revisao}

\subsection{Geração aumentada por recuperação}

Na formulação de Lewis et al.~\cite{lewis2020rag}, um sistema RAG combina um componente de recuperação, que seleciona trechos de uma coleção de documentos a partir da pergunta, e um componente gerador, que condiciona a resposta a esses trechos. Em implementações atuais, os documentos são divididos em fragmentos (\emph{chunks}) \cite{schwaber2025chunking}, convertidos em vetores por um modelo de embedding e armazenados em um banco vetorial \cite{ma2024vectorsurvey}; a pergunta é codificada pelo mesmo modelo, e os $k$ fragmentos mais próximos são entregues ao LLM. A arquitetura reduz alucinações e permite atualizar o conhecimento sem retreinar o modelo \cite{andriopoulos2023augmenting}, mas sua robustez depende de decisões tomadas empiricamente: Barnett et al.~\cite{barnett2024seven} catalogam sete pontos de falha recorrentes, entre eles a recuperação de trechos errados e a incapacidade do gerador de extrair a resposta de um contexto que a contém.

\subsection{Métodos de injeção de contexto}

Uma vez recuperados os fragmentos, há diferentes formas de entregá-los ao LLM \cite{topsakal2023langchain}. No \textbf{Stuff}, todos os fragmentos são concatenados em um único prompt. No \textbf{Refine}, a resposta é construída a partir do primeiro fragmento e refinada sequencialmente a cada fragmento seguinte. No \textbf{Map-Reduce}, cada fragmento gera uma resposta parcial independente, e uma chamada final combina as parciais. No \textbf{Map-Rerank}, cada fragmento gera uma resposta acompanhada de uma nota de confiança, e mantém-se apenas a de maior confiança. Şakar e Emekci~\cite{sakar2025rag} acrescentam dois métodos que atuam sobre a pergunta: o \textbf{Query Step-Down} gera perguntas alternativas mais específicas, responde cada uma e combina as respostas; o \textbf{Reciprocal RAG} segue o mesmo primeiro passo, mas retém apenas a resposta mais confiável. Ambos dialogam com a reescrita de consultas \cite{ma2023query} e com o RAG-Fusion \cite{rackauckas2024ragfusion}, que ampliam a cobertura da recuperação ao custo de chamadas adicionais ao LLM e do risco de desvio em relação à intenção original da pergunta.

A escolha do método tem efeito direto sobre o número de chamadas ao LLM e, portanto, sobre tokens e latência. Há também um efeito sobre a qualidade: Liu et al.~\cite{liu2023lost} mostram que LLMs aproveitam pior a informação posicionada no meio de contextos longos, o que penaliza estratégias que apenas acumulam fragmentos, e Nair et al.~\cite{nair2023drilling} mostram que representações condensadas de documentos longos preservam quase todo o desempenho com uma fração dos tokens.

\subsection{Bancos vetoriais e busca aproximada}

Bancos vetoriais como o ChromaDB oferecem persistência em disco, armazenamento dos textos e filtros por metadados, e, por padrão, indexam os vetores com o HNSW \cite{malkov2020hnsw}, um grafo de proximidade em camadas que encontra vizinhos próximos em tempo sublinear, mas de forma aproximada: os $k$ fragmentos devolvidos podem não ser os $k$ mais próximos. Bibliotecas como o FAISS \cite{johnson2021faiss} oferecem tanto índices aproximados quanto a busca exata por força bruta (\texttt{IndexFlatL2}), cujo custo cresce linearmente com o número de vetores. A diferença entre os dois bancos, portanto, não é apenas de infraestrutura: com o índice aproximado, a escolha do banco pode mudar os fragmentos entregues ao LLM, efeito que tende a crescer com o tamanho da base.

\subsection{Avaliação de sistemas RAG}

A similaridade de cosseno entre os embeddings da resposta gerada e da resposta de referência é uma métrica barata e reprodutível, mas mede proximidade semântica global: penaliza respostas corretas redigidas de forma mais longa que a referência, favorece respostas no mesmo vocabulário mesmo quando o conteúdo está errado e não verifica se a resposta está ancorada nos documentos recuperados. Em conjuntos de compreensão de leitura, a métrica usual é a sobreposição de tokens com a resposta de referência (F1 e \emph{exact match}) \cite{rajpurkar2016squad}, e, em questões de múltipla escolha, a acurácia.

O \emph{framework} RAGAS \cite{es2024ragas} propõe métricas específicas para RAG calculadas com um LLM como juiz: \emph{Faithfulness} (fração das afirmações da resposta sustentadas pelo contexto recuperado), \emph{Answer Relevancy} (quanto a resposta endereça a pergunta), \emph{Context Precision} e \emph{Context Recall} (qualidade e suficiência do contexto recuperado em relação à referência). O uso de LLMs como juízes, porém, tem vieses documentados, entre eles o de autopreferência: juízes tendem a avaliar melhor textos gerados por eles mesmos \cite{zheng2023judging, panickssery2024selfpref}. O RAGBench \cite{friel2024ragbench} oferece respostas anotadas quanto à sustentação pelo contexto em sete subconjuntos de domínios variados, o que permite medir a concordância de um juiz automático com um gabarito.

\subsection{O papel do LLM e do domínio}

O ganho do RAG não é uniforme. Mallen et al.~\cite{mallen2023trust} mostram que a recuperação ajuda sobretudo em conhecimento pouco popular, que o modelo não memorizou, e pode até prejudicar em fatos populares, que o modelo já responde de memória. No domínio médico, o MIRAGE \cite{xiong2024mirage} avaliou o RAG com vários LLMs e corpora, entre eles os livros-texto que acompanham o MedQA, sempre contra o mesmo modelo sem recuperação, e encontrou ganhos que variam com o modelo e com o corpus. A variação do LLM também importa por razões de custo e de reprodutibilidade: o FrugalGPT \cite{chen2023frugalgpt} mostra que combinações de modelos mais baratos podem alcançar o desempenho de modelos caros a uma fração do custo; Chen, Zaharia e Zou~\cite{chen2024chatgptbehavior} documentam mudanças de comportamento de um mesmo modelo comercial ao longo de poucos meses; e Palmer, Smith e Spirling~\cite{palmer2024proprietary} defendem que pesquisas com modelos proprietários justifiquem a escolha e incluam modelos abertos, cujos pesos permitem reproduzir os resultados.

\subsection{RAG em português}

A avaliação de RAG em português ainda é incipiente. Medeiros e Oliveira~\cite{medeiros2025embeddings} compararam modelos de embedding e LLMs abertos e proprietários em três bases em português (Pirá, SQuAD 2 e FairytaleQA PT-BR), com ROUGE-L, BERTScore e as métricas RAGAS de fidelidade e relevância da resposta. Os melhores resultados foram obtidos pelo Multilingual E5 large e pelo Gemma 2 9B. Aquele estudo varia o embedding e o LLM com um método de injeção fixo; o presente trabalho faz o recorte complementar, fixando o embedding e variando o método de injeção e o LLM.

\subsection{Lacuna}

A literatura avançou na acurácia \cite{ma2023query, rackauckas2024ragfusion}, na confiabilidade \cite{andriopoulos2023augmenting, barnett2024seven} e no entendimento das limitações de contexto \cite{liu2023lost, nair2023drilling} de sistemas RAG. A comparação sistemática entre métodos de injeção de contexto, porém, foi feita com documentos em inglês, uma única métrica de similaridade e sem um controle sem recuperação \cite{sakar2025rag}, enquanto os estudos que variam o LLM e incluem esse controle fixam o método de injeção \cite{xiong2024mirage, medeiros2025embeddings}. Não há, até onde foi possível verificar, uma comparação dos métodos de injeção que cruze domínios, idiomas e LLMs com um controle sem RAG e métricas verificadas contra o gabarito. Este trabalho se insere nessa lacuna.
